\documentclass[11pt,aspectratio=169]{beamer}

\usetheme{Madrid}
\usecolortheme{seahorse}

\usepackage[utf8]{inputenc}
\usepackage{amsmath,amssymb}
\usepackage{xcolor}
\usepackage{booktabs}
\usepackage{tikz}
\usetikzlibrary{arrows.meta,positioning,fit,backgrounds}

\setbeamertemplate{navigation symbols}{}
\setbeamertemplate{footline}[frame number]

\title{Beyond the Syllabus}
\subtitle{How Fast Can Matching Go? Continuous Methods, Combinatorial Counterattacks,
and Online Arrivals\\
\small Week 11 Supplement -- TOAE209/TOAH209: Design and Analysis of Algorithms}
\author{Foreign Trade University}
\date{}

\begin{document}

% ---------------------------------------------------------------
\begin{frame}
  \titlepage
\end{frame}
% ---------------------------------------------------------------

\begin{frame}{Outline}
  \tableofcontents
\end{frame}

% =================================================================
\section{Motivation}
% =================================================================

\begin{frame}{Why look beyond the syllabus?}
  \small
  \begin{itemize}
    \item This week you reduced \textbf{maximum bipartite matching} to a max-flow
    computation: build $s\to X \to Y \to t$ with unit capacities, run Edmonds--Karp,
    read off the matching. Since Edmonds--Karp augments along shortest paths, this
    reduction runs in $O(VE)$ -- and for bipartite graphs specifically, the classic
    Hopcroft--Karp algorithm (a smarter augmenting-path schedule) gets this down to
    $O(E\sqrt{V})$, a bound that stood as the best known for almost \textbf{fifty years}
    (1973--2020).
    \item Three questions a curious student should ask next:
    \begin{enumerate}
      \item Can matching beat $O(E\sqrt{V})$ at all -- and if so, how?
      \item Does beating it \emph{require} heavy continuous-optimization machinery, or
      can pure combinatorics get there too?
      \item Everything above assumes the whole graph is known upfront. What if vertices
      of $Y$ (jobs, ads, riders) \textbf{arrive online}, one at a time, and you must
      match them irrevocably as they come?
    \end{enumerate}
  \end{itemize}
  \begin{alertblock}{How this supplement is different from a survey}
    Every section is grounded in a specific, verifiable paper (with arXiv/DOI); none of
    the bounds below are approximate folklore.
  \end{alertblock}
\end{frame}

% =================================================================
\section{Question 1: Breaking the 1973 Barrier}
% =================================================================

\begin{frame}{Fifty years of $O(E\sqrt{V})$}
  \begin{itemize}
    \item Hopcroft \& Karp (1973) showed that after $O(\sqrt{V})$ phases of
    \emph{shortest-augmenting-path} search, a bipartite matching is maximum; each
    phase costs $O(E)$, giving $O(E\sqrt{V})$ total.
    \item Dinic (1970) and Karzanov (1973) proved essentially the same bound
    independently, via blocking flows -- three different derivations of the same wall.
    \item For dense graphs ($E = \Theta(V^2)$), $O(E\sqrt{V}) = O(V^{2.5})$: at
    $V=10^5$, that is $\sim 3\times10^{12}$ operations -- far too slow.
  \end{itemize}
\end{frame}

\begin{frame}{The break: continuous optimization enters combinatorics}
  \begin{block}{van den Brand, Lee, Nanongkai, Peng, Saranurak, Sidford, Song \& Wang
  (FOCS 2020)}
    \emph{``Bipartite Matching in Nearly-linear Time on Moderately Dense Graphs.''}
    arXiv:2009.01802. A randomized algorithm solving maximum bipartite matching (and
    several relatives: transshipment, negative-weight shortest paths, optimal
    transport) in $\tilde{O}(m + n^{1.5})$ time.
  \end{block}
  \begin{itemize}
    \item For moderately dense graphs ($m = \Omega(n^{1.5})$), this is
    \textbf{nearly linear in the input size} -- the first improvement over
    Dinic/Hopcroft--Karp/Karzanov's $O(m\sqrt{n})$ in 47 years.
    \item \textbf{How:} recast matching as a linear program, solve it with an
    \textbf{interior-point method} (IPM). Each IPM iteration reduces to solving a
    \textbf{Laplacian linear system} on the graph -- and those have had
    nearly-linear-time solvers since Spielman \& Teng (2004). Dynamic graph data
    structures avoid recomputing each Laplacian solve from scratch as the IPM
    progresses.
  \end{itemize}
\end{frame}

\begin{frame}{Putting numbers on it}
  \small
  At $n = 10^6$, $m = n^{1.5} \approx 10^9$:
  \begin{center}
  \begin{tabular}{@{}lll@{}}
    \toprule
    Algorithm & Bound & Rough op.\ count \\
    \midrule
    Hopcroft--Karp / Dinic (1970s) & $O(m\sqrt{n})$ & $\sim 10^{12}$ \\
    van den Brand et al.\ (FOCS 2020) & $\tilde{O}(m + n^{1.5})$ & $\sim 10^9$ \\
    \bottomrule
  \end{tabular}
  \end{center}
  Three orders of magnitude, from a fifty-year-old wall, by importing a tool
  (Laplacian solvers) that this week's lecture never mentions -- matching is not just
  a graph problem once you let continuous optimization in.
\end{frame}

% =================================================================
\section{Question 2: Do You Actually Need All That Machinery?}
% =================================================================

\begin{frame}{A combinatorial counterattack}
  \begin{block}{Chuzhoy \& Khanna (SODA 2024 / STOC 2024)}
    \emph{``A Faster Combinatorial Algorithm for Maximum Bipartite Matching''}
    (arXiv:2312.12584, SODA 2024, $\tilde O(m^{1/3}n^{5/3})$-time) and
    \emph{``Maximum Bipartite Matching in $n^{2+o(1)}$ Time via a Combinatorial
    Algorithm''} (arXiv:2405.20861, STOC 2024).
  \end{block}
  \begin{itemize}
    \item Both results match or beat the IPM-based bound above -- using \textbf{no}
    interior-point methods, \textbf{no} Laplacian solvers, only combinatorial
    augmenting-path ideas pushed further than Hopcroft--Karp's.
    \item The lesson: the 2020 breakthrough proved \emph{that} $O(m\sqrt n)$ could be
    beaten, but it did \textbf{not} prove that continuous machinery was
    \emph{necessary}. That is a separate, harder question -- and as of 2024 the answer
    leans ``no, elbow grease still works too.''
  \end{itemize}
  \begin{exampleblock}{Why this matters pedagogically}
    Two independent, correct routes to the same destination is exactly the kind of
    thing worth noticing: a faster \emph{bound} does not tell you the \emph{only} way
    to achieve it.
  \end{exampleblock}
\end{frame}

% =================================================================
\section{Question 3: What If the Graph Arrives Online?}
% =================================================================

\begin{frame}{The online model: Karp--Vazirani--Vazirani (1990)}
  \begin{itemize}
    \item Everything above assumes $X$, $Y$, and all edges are known before you start.
    In \textbf{online bipartite matching}, $X$ is fixed but vertices of $Y$ arrive one
    at a time, each revealing its edges to $X$ on arrival; you must match it
    immediately (or leave it unmatched) with \textbf{no take-backs}.
    \item Karp, Vazirani \& Vazirani (STOC 1990) proved the \textsc{Ranking}
    algorithm -- pick one random permutation of $X$ up front, always match an
    arriving $y$ to its \emph{earliest-ranked} available neighbor -- achieves
    competitive ratio $1 - 1/e \approx 0.632$ against the offline optimum, and that no
    online algorithm can beat $1-1/e$ in the worst case. A genuinely tight bound,
    open for `` how do we do better'' ever since -- \emph{if} the adversary is fully
    unconstrained.
  \end{itemize}
\end{frame}

\begin{frame}{Beating $1-1/e$ with (untrusted) predictions}
  \begin{block}{Choo, Gouleakis, Ling \& Bhattacharyya (ICML 2024)}
    \emph{``Online bipartite matching with imperfect advice.''} PMLR 235:8762--8781,
    arXiv:2405.09784.
  \end{block}
  \begin{itemize}
    \item Setting: before the online phase, the algorithm is handed a
    \textbf{prediction} of the arrival sequence -- possibly wrong, possibly
    adversarially wrong.
    \item Idea: use a \textbf{prefix} of the actual arrivals as a statistical test of
    whether the prediction looks trustworthy; if it does, follow it (and beat
    $1-1/e$ when it is accurate); if not, fall back to a \textsc{Ranking}-style
    worst-case-safe baseline.
    \item This is the \textbf{algorithms-with-predictions} paradigm (Mitzenmacher \&
    Vassilvitskii): a \emph{consistency} guarantee (near-optimal if predictions are
    good) traded against a \emph{robustness} guarantee (never worse than $1-1/e$-ish
    even if predictions are adversarial) -- both guarantees proven, not just
    empirically observed.
  \end{itemize}
\end{frame}

% =================================================================
\section{Summary}
% =================================================================

\begin{frame}{Summary}
  \small
  \begin{enumerate}
    \item \textbf{The 1973 barrier fell in 2020:} van den Brand et al.\ (FOCS 2020)
    solve bipartite matching in $\tilde O(m+n^{1.5})$ by routing through Laplacian
    solvers and interior-point methods -- continuous mathematics inside a purely
    combinatorial problem.
    \item \textbf{Machinery is not mandatory:} Chuzhoy \& Khanna (SODA/STOC 2024) reach
    a comparable bound with a purely combinatorial algorithm -- there is more than one
    road to a faster algorithm.
    \item \textbf{``Known upfront'' is a strong assumption:} once arrivals are online,
    even the optimal worst-case algorithm (Karp--Vazirani--Vazirani, 1990) is stuck at
    $1-1/e$ -- and only a genuinely new resource, untrusted \emph{predictions}
    (Choo et al., ICML 2024), lets you provably beat that wall while staying safe when
    the predictions are wrong.
  \end{enumerate}
  \begin{alertblock}{Looking ahead}
    Next week (Kleinberg \& Tardos, Ch.\ 8) asks \emph{which problems provably have no
    polynomial-time exact algorithm} -- NP-completeness. ``Optimal Account Balancing,''
    this week's Hard problem, is already a preview: minimizing the transaction
    \emph{count} has no known polynomial algorithm.
  \end{alertblock}
\end{frame}

\end{document}
